\documentclass[tikz,border=7pt]{standalone}
\usepackage[american]{circuitikz}
\standaloneenv{circuitikz}
\begin{document}
\begin{circuitikz}[font=\sffamily]
 \draw (-3.6,3.2) -- (3.6,3.2) node[right] {$+V_{CC}$};
 \draw (-2,2.15) node[pnp,xscale=-1](Q3){};
 \node at (-2.55,2.15) {$Q_3$};
 \draw (2,2.15) node[pnp](Q4){$Q_4$};
 \draw (Q3.E) -- (-2,3.2); \draw (Q4.E) -- (2,3.2);
 \draw (Q3.B) -- (Q4.B);
 \draw (Q3.C) -- (-2,.5) node[circ](L){};
 \draw (Q3.B) -- (-.65,2.15) -- (-.65,.5) -- (L);
 \draw (-2,.5) -- (-2,.15) node[npn,anchor=C](Q1){$Q_1$};
 \draw (2,.5) node[circ,label=right:$v_o$]{} -- (2,.15)
       node[npn,anchor=C,xscale=-1](Q2){};
 \node at (2.55,-1.05) {$Q_2$};
 \draw (Q1.B) -- (-3.4,-.62) node[left] {$v_1$};
 \draw (Q2.B) -- (3.4,-.62) node[right] {$v_2$};
 \draw (Q1.E) -- (-2,-1.65) -- (2,-1.65) -- (Q2.E);
 \draw (0,-1.65) to[I,l_=$I_T$] (0,-3) node[ground]{};
 \draw (Q4.C) -- (2,.5);
 \draw[->,teal!75!black,very thick] (-1.3,.3) .. controls (-.3,1) and (.3,1) .. (1.3,.3)
   node[midway,above] {mirrored signal current};
 \node[align=center,fill=white] at (0,-3.65)
 {PNP mirror active load converts the two collector-current changes\\
 into one approximately doubled single-ended signal current.};
\end{circuitikz}
\end{document}
